% Section 6: what the report established.
\section{Conclusion}
\label{sec:conclusion}

The requirements force the maze to be a spanning tree of the grid, and
recursive division builds one in $\Theta(n)$ while recording, at no extra cost,
the binary tree of divisions that made it.

That second structure changes what solving costs. On the graph, A* returns the
unique route in $O(n\log n)$ and in practice examines most of the maze, because
a tree offers a heuristic nothing to prune. On the binary tree, a lowest common
ancestor finds the first wall the route crosses in $\Theta(\log n)$, and the
whole route in $\Theta(n^{0.68})$ --- thirty times fewer visits at
$256\times256$, with logarithmic memory.

The generator knows where the walls are. Keeping that knowledge, rather than
rediscovering it by search, is the whole of the gain.
